\section{Market regimes}
\label{sec:regimes}

Statistical-arbitrage returns are known to depend on the market
environment. Reversion profits are typically higher when dispersion is
high, and they can collapse in crowded deleveraging episodes
\citep{khandani2011}. We summarize the environment with three daily
features computed from data through day~$d$: the log 21-day realized
volatility of the equal-weight market return; the 63-day average
pairwise correlation; and the log 21-day average cross-sectional return
dispersion. The average correlation is obtained without forming any
matrix. For $N$ unit-variance assets with average correlation
$\bar\rho$, the cross-sectional mean of volatility-standardized returns
has variance $\bar\rho + (1-\bar\rho)/N$, which we invert.

A two-state Gaussian hidden Markov model \citep{hamilton1989} with
diagonal covariances is fit by Baum--Welch \citep{rabiner1989}. To
avoid look-ahead, the model is refit once a year on all data before
1~January, with features standardized using that training set's
moments. The probabilities passed downstream are \emph{filtered}
(forward-algorithm) probabilities $P(\text{stress}\mid x_{1:d})$, never
smoothed ones. States are ordered by mean volatility, so state~1 is
always the high-volatility state, and the stress probability on the
day before an episode starts enters both the ML models (as a feature)
and the portfolio (as an exposure overlay).
